\chapter{绪论}

Kresling 折纸单元兼具几何可展开性与轴向变形能力，可用于构造轻量化、可重构的工程结构。本章示例用于展示硕士论文正文的章节标题、页眉、页码及引用格式。正式写作时，应补充研究背景、国内外研究现状、当前问题以及本文的主要工作，并规范引用文献，例如 \cite{miura1985method,kresling_example}。

\section{研究背景与意义}

示例文字用于检查中英文混排、单栏宽度和段落间距。替换内容时建议保留语义完整的段落，并使用 \verb|\label| 与 \verb|\cref| 管理图、表和公式交叉引用。

\section{国内外研究现状}

本节可按几何建模、力学响应、驱动方式和工程应用等主题综述已有研究，明确本文工作的切入点。

\section{本文主要研究内容}

本文后续章节依次给出理论模型、数值与实验方法、结果分析以及结论与展望。
